% GENERATED FILE. Edit canonical lessons, then run npm run generate:books.
% canonical-source-hash: fnv1a64:c4d894404c712cd9
% canonical-lessons: BN-C402-majhkhan, BN-C402-ghurpoth, BN-C402-kinara, BN-C402-dishehara, BN-C402-ghurpak-khaoya

\chapter{Middles and Edges}
\label{ch:bn-middles-and-edges}
\hlchaptermodality{402}

\begin{chapteropening}
\textbf{By the end of this chapter:} \emph{I can say where something stands in the middle or at the edge, and that we went the long way round.}

Complete the last lesson of chapter 402: I can say where something stands in the middle or at the edge, and that we went the long way round.
\end{chapteropening}

\section[\texorpdfstring{\bn{মাঝখান} (mājhkhān)}{মাঝখান (mājhkhān)}]{\bn{মাঝখান} (mājhkhān) — the middle}

\label{lesson:BN-C402-majhkhan}



\begin{quote}
Before the new one: say the Bengali for a compass, then the Bengali for the horizon.
\end{quote}

\subsection*{You'll want to know: \bn{মাঝখান}}
\textbf{\bn{মাঝখান}} — \emph{mājhkhān} — \textquotedblleft{}the middle\textquotedblright{}.

\begin{quote}
\bn{রাস্তার} \bn{মাঝখানে} \bn{একটা} \bn{গরু} \bn{দাঁড়িয়ে} \bn{ছিল।}
\end{quote}

\begin{quote}
\emph{rāstār mājhkhāne ekṭā gôru dā̃ṛiye chhilô.}
\end{quote}

\begin{quote}
A cow was standing in the middle of the road.
\end{quote}

\bn{মাঝখানে}, in the middle; the \bn{মাঝ} of \bn{মাঝে} \bn{মাঝে}, sometimes, is the same word.

\begin{cousinweb}
\textbf{\bn{মাঝখান}} (\emph{mājhkhān}) joins \bn{মাঝ}, \textquotedblleft{}the middle\textquotedblright{}, which CDIAL lists as Bengali \emph{mājh} under Sanskrit \emph{mádhya}, \textquotedblleft{}middle\textquotedblright{} (CDIAL 9804), to \bn{খান}, \textquotedblleft{}a place\textquotedblright{}, the ending of \bn{কারখানা}.
\end{cousinweb}

\begin{grammarlens}[title={What you've built}]
One more word for getting around.
\end{grammarlens}

\subsection*{Your turn}
\begin{itemize}
\raggedright
  \item \emph{Say it:} \emph{mājhkhān}
  \item \emph{Say it:} \emph{mājhkhān}, once more
  \item \emph{Say it:} \emph{rāstār mājhkhāne ekṭā gôru dā̃ṛiye chhilô.} — A cow was standing in the middle of the road.
  \item \emph{Recall:} say the Bengali for a compass, then the Bengali for the horizon, then say \emph{mājhkhān} again
\end{itemize}

\begin{tcolorbox}[breakable,colback=teal!4,colframe=teal!35!black,title={Before you move on}]
What does \bn{মাঝখান} mean? (\textquotedblleft{}The middle\textquotedblright{}.) Say it once more.
\end{tcolorbox}
\section[\texorpdfstring{\bn{ঘুরপথ} (ghurpôth)}{ঘুরপথ (ghurpôth)}]{\bn{ঘুরপথ} (ghurpôth) — a roundabout way; a detour}

\label{lesson:BN-C402-ghurpoth}



\begin{quote}
Before the new one: say the Bengali for the horizon, then the Bengali for the middle.
\end{quote}

\subsection*{You'll want to know: \bn{ঘুরপথ}}
\textbf{\bn{ঘুরপথ}} — \emph{ghurpôth} — \textquotedblleft{}a roundabout way; a detour\textquotedblright{}.

\begin{quote}
\bn{সেতু} \bn{বন্ধ} \bn{ছিল}, \bn{তাই} \bn{আমরা} \bn{ঘুরপথে} \bn{গেলাম।}
\end{quote}

\begin{quote}
\emph{setu bôndhô chhilô, tāi āmrā ghurpôthe gelām.}
\end{quote}

\begin{quote}
The bridge was closed, so we went the long way round.
\end{quote}

\begin{cousinweb}
\textbf{\bn{ঘুরপথ}} (\emph{ghurpôth}) joins the stem \bn{ঘুর}- of \bn{ঘোরা}, \textquotedblleft{}to go round\textquotedblright{} (CDIAL 4488), to \textbf{\bn{পথ}}, \textquotedblleft{}a way\textquotedblright{}: the way that goes round.
\end{cousinweb}

\begin{grammarlens}[title={What you've built}]
One more word for getting around.
\end{grammarlens}

\subsection*{Your turn}
\begin{itemize}
\raggedright
  \item \emph{Say it:} \emph{ghurpôth}
  \item \emph{Say it:} \emph{ghurpôth}, once more
  \item \emph{Say it:} \emph{setu bôndhô chhilô, tāi āmrā ghurpôthe gelām.} — The bridge was closed, so we went the long way round.
  \item \emph{Recall:} say the Bengali for the horizon, then the Bengali for the middle, then say \emph{ghurpôth} again
\end{itemize}

\begin{tcolorbox}[breakable,colback=teal!4,colframe=teal!35!black,title={Before you move on}]
What does \bn{ঘুরপথ} mean? (\textquotedblleft{}A roundabout way; a detour\textquotedblright{}.) Say it once more.
\end{tcolorbox}
\section[\texorpdfstring{\bn{কিনারা} (kinārā)}{কিনারা (kinārā)}]{\bn{কিনারা} (kinārā) — an edge; a brink; a shore}

\label{lesson:BN-C402-kinara}



\begin{quote}
Before the new one: say the Bengali for the middle, then the Bengali for a detour.
\end{quote}

\subsection*{You'll want to know: \bn{কিনারা}}
\textbf{\bn{কিনারা}} — \emph{kinārā} — \textquotedblleft{}an edge; a brink; a shore\textquotedblright{}.

\begin{quote}
\bn{অরুণ} \bn{খাদের} \bn{কিনারা} \bn{থেকে} \bn{নিচে} \bn{তাকাল।}
\end{quote}

\begin{quote}
\emph{ôrun khāder kinārā theke niche tākālô.}
\end{quote}

\begin{quote}
Arun looked down from the edge of the ravine.
\end{quote}

Also a way out of a problem: \bn{কিনারা} \bn{করা}, to find a solution.

\begin{cousinweb}
\textbf{\bn{কিনারা}} (\emph{kinārā}) is usually traced to Persian \emph{kinār}, \textquotedblleft{}an edge, a bank, a shore\textquotedblright{}.
\end{cousinweb}

\begin{grammarlens}[title={What you've built}]
One more word for getting around.
\end{grammarlens}

\subsection*{Your turn}
\begin{itemize}
\raggedright
  \item \emph{Say it:} \emph{kinārā}
  \item \emph{Say it:} \emph{kinārā}, once more
  \item \emph{Say it:} \emph{ôrun khāder kinārā theke niche tākālô.} — Arun looked down from the edge of the ravine.
  \item \emph{Recall:} say the Bengali for the middle, then the Bengali for a detour, then say \emph{kinārā} again
\end{itemize}

\begin{tcolorbox}[breakable,colback=teal!4,colframe=teal!35!black,title={Before you move on}]
What does \bn{কিনারা} mean? (\textquotedblleft{}An edge; a brink; a shore\textquotedblright{}.) Say it once more.
\end{tcolorbox}
\section[\texorpdfstring{\bn{দিশেহারা} (dishehārā)}{দিশেহারা (dishehārā)}]{\bn{দিশেহারা} (dishehārā) — at a loss; not knowing which way to go}

\label{lesson:BN-C402-dishehara}



\begin{quote}
Before the new one: say the Bengali for a detour, then the Bengali for an edge.
\end{quote}

\subsection*{You'll want to know: \bn{দিশেহারা}}
\textbf{\bn{দিশেহারা}} — \emph{dishehārā} — \textquotedblleft{}at a loss; not knowing which way to go\textquotedblright{}.

\begin{quote}
\bn{নতুন} \bn{শহরে} \bn{এসে} \bn{আমি} \bn{দিশেহারা} \bn{হয়ে} \bn{গেলাম।}
\end{quote}

\begin{quote}
\emph{nôtun shôhôre ese āmi dishehārā hôye gelām.}
\end{quote}

\begin{quote}
When I came to the new city I was at a loss.
\end{quote}

\begin{cousinweb}
\textbf{\bn{দিশেহারা}} (\emph{dishehārā}) joins \bn{দিশে}, a spoken form of \bn{দিশা}, \textquotedblleft{}direction\textquotedblright{}, Sanskrit \emph{diśā}, \textquotedblleft{}a direction, a quarter\textquotedblright{} (Monier-Williams), to \bn{হারা}, \textquotedblleft{}having lost\textquotedblright{}, the stem of \bn{হারানো}: one who has lost the direction.
\end{cousinweb}

\begin{grammarlens}[title={What you've built}]
One more word for getting around.
\end{grammarlens}

\subsection*{Your turn}
\begin{itemize}
\raggedright
  \item \emph{Say it:} \emph{dishehārā}
  \item \emph{Say it:} \emph{dishehārā}, once more
  \item \emph{Say it:} \emph{nôtun shôhôre ese āmi dishehārā hôye gelām.} — When I came to the new city I was at a loss.
  \item \emph{Recall:} say the Bengali for a detour, then the Bengali for an edge, then say \emph{dishehārā} again
\end{itemize}

\begin{tcolorbox}[breakable,colback=teal!4,colframe=teal!35!black,title={Before you move on}]
What does \bn{দিশেহারা} mean? (\textquotedblleft{}At a loss; not knowing which way to go\textquotedblright{}.) Say it once more.
\end{tcolorbox}
\section[\texorpdfstring{\bn{ঘুরপাক} \bn{খাওয়া} (ghurpāk khāoyā)}{ঘুরপাক খাওয়া (ghurpāk khāoyā)}]{\bn{ঘুরপাক} \bn{খাওয়া} (ghurpāk khāoyā) — to go round in circles; to whirl}

\label{lesson:BN-C402-ghurpak-khaoya}



\begin{quote}
Before the new one: say the Bengali for an edge, then the Bengali for at a loss.
\end{quote}

\subsection*{You'll want to know: \bn{ঘুরপাক} \bn{খাওয়া}}
\textbf{\bn{ঘুরপাক} \bn{খাওয়া}} — \emph{ghurpāk khāoyā} — \textquotedblleft{}to go round in circles; to whirl\textquotedblright{}.

\begin{quote}
\bn{পথ} \bn{না} \bn{জেনে} \bn{আমরা} \bn{এক} \bn{ঘণ্টা} \bn{ঘুরপাক} \bn{খেলাম।}
\end{quote}

\begin{quote}
\emph{pôth nā jene āmrā ek ghônṭā ghurpāk khelām.}
\end{quote}

\begin{quote}
Not knowing the way, we went round in circles for an hour.
\end{quote}

Past (he/she …): \bn{তিনি} \bn{ঘুরপাক} \bn{খেলেন}, \emph{tini ghurpāk khelen}.

\bn{খাওয়া}, to eat, also means to undergo: \bn{ধাক্কা} \bn{খাওয়া}, to take a knock; \bn{ঘুরপাক} \bn{খাওয়া}, to take turn after turn.

\begin{cousinweb}
\textbf{\bn{ঘুরপাক} \bn{খাওয়া}} (\emph{ghurpāk khāoyā}) joins \bn{ঘুর}-, \textquotedblleft{}going round\textquotedblright{} (CDIAL 4488), and \bn{পাক}, which CDIAL lists as Bengali \emph{pāk}, \textquotedblleft{}a bend, a twist, an eddy\textquotedblright{} (CDIAL 7618), to \bn{খাওয়া}, to eat, here to undergo.
\end{cousinweb}

\begin{grammarlens}[title={What you've built}]
One more word for getting around.
\end{grammarlens}

\subsection*{Your turn}
\begin{itemize}
\raggedright
  \item \emph{Say it:} \emph{ghurpāk khāoyā}
  \item \emph{Say it:} \emph{ghurpāk khāoyā}, once more
  \item \emph{Say it:} \emph{pôth nā jene āmrā ek ghônṭā ghurpāk khelām.} — Not knowing the way, we went round in circles for an hour.
  \item \emph{Recall:} say the Bengali for an edge, then the Bengali for at a loss, then say \emph{ghurpāk khāoyā} again
\end{itemize}

\begin{tcolorbox}[breakable,colback=teal!4,colframe=teal!35!black,title={Before you move on}]
What does \bn{ঘুরপাক} \bn{খাওয়া} mean? (\textquotedblleft{}To go round in circles; to whirl\textquotedblright{}.) Say it once more.
\end{tcolorbox}
